\section{Related Work}

\subsection{Polar AWS Records and Reanalysis Products}
Antarctic AWS records are indispensable for evaluating surface meteorology in regions where staffed stations are sparse. AWS networks provide the only sustained in situ observations for many coastal, ice-shelf, and interior locations, but their records are also shaped by the same environmental constraints that make them valuable: winter darkness, icing, snow burial, power interruptions, and infrequent maintenance access. The AntAWS dataset compiles quality-controlled AWS observations across Antarctica and provides 3-hourly meteorological variables including air temperature, pressure, relative humidity, and wind \citep{wang2023antaws}. The IMAU Antarctic AWS archive adds a complementary long-term record with additional surface-energy and surface-height variables at 19 locations from 1995--2022 \citep{vantiggelen2025imau}. Although Greenland is not the target region of this work, the PROMICE and GC-NET AWS archives are important methodological comparators because they document the kind of sensor metadata, quality-control logic, station redesign, and operational data release practices expected for polar station datasets \citep{fausto2021promice,fausto2026promicegcnet}.

Reanalysis products provide spatially complete, physically constrained background fields that are attractive for sparse polar settings. ERA5 provides hourly global atmospheric variables and is widely used as a modern reanalysis baseline \citep{hersbach2020era5}. However, Antarctic validation studies show that reanalysis should not be treated as ground truth. ERA5 near-surface temperature over Antarctica has high monthly correlation with station observations but exhibits regionally structured biases \citep{zhu2021era5antarctica}. In the Antarctic Peninsula--Ellsworth Land region, ERA5 performance depends on parameter and geographic context \citep{tetzner2019era5antarcticpeninsula}. Coastal wind studies likewise show that ERA5 can be the strongest available reanalysis while still missing local variability near complex orography and coastal margins \citep{catonharrison2022antarcticwinds}. These findings support the design choice in ECBIT: ERA5 should be used as a conditional prior during outages, not as a wholesale replacement for AWS observations.

Recent climate-data reconstruction work further motivates combining sparse observations with physically coherent gridded products. \citet{ma2025antarcticsat} train a deep model with reanalysis fields and observation masks to reconstruct Antarctic surface air temperature from sparse in situ records. \citet{li2026era5superresolution} study super-resolution transferability for ERA5, highlighting the broader role of machine learning in refining reanalysis information across regions and variables. ECBIT differs from these gridded reconstruction and downscaling tasks because its target is station-level temporal imputation inside sparse AWS windows, where observed AWS values must be preserved and ERA5 should be activated only for missing variables and timesteps.

\subsection{Time-Series Imputation}
Classical imputation baselines such as linear interpolation and last-observation-carried-forward are strong sanity checks for short gaps but do not model cross-variable structure. Non-parametric methods such as MissForest can exploit cross-sectional feature relations but are not designed around ordered temporal dynamics \citep{stekhoven2012missforest}. Deep imputation methods add temporal modeling and learned latent structure. GRU-D incorporates missingness indicators and time gaps into recurrent modeling \citep{che2018grud}. BRITS learns missing values inside a bidirectional recurrent computation graph \citep{cao2018brits}. SAITS uses self-attention and missingness-aware weighting for multivariate time-series imputation \citep{du2023saits}. GP-VAE and CSDI frame imputation probabilistically through latent-variable or score-based generative modeling \citep{fortuin2020gpvae,tashiro2021csdi}.

Survey and benchmark work makes clear that time-series imputation performance is tightly coupled to missingness assumptions. \citet{fang2020imputationsurvey} review deep imputation architectures and their treatment of temporal structure, while \citet{wang2024imputationsurvey} organize multivariate imputation methods by uncertainty modeling and neural architecture. TSI-Bench standardizes deep imputation evaluation across many algorithms, datasets, and missingness scenarios \citep{du2024tsibench}. Its scale is valuable, but ECBIT addresses a different question. TSI-Bench asks how generic imputation algorithms compare under benchmark missingness protocols; ECBIT asks whether a domain-specific block-missing curriculum and ERA5 conditioning improve Antarctic AWS recovery under realistic contiguous station outages. The protocols are therefore complementary rather than interchangeable.

\subsection{Transformers for Time Series}
ECBIT uses transformer encoders because self-attention provides a flexible mechanism for variable-level interaction and long-window context modeling \citep{vaswani2017attention}. Time-series transformers have evolved away from direct transplantation of language-model tokenization. Informer and Autoformer introduced efficient attention and decomposition ideas for long-sequence forecasting \citep{zhou2021informer,wu2021autoformer}. FEDformer combines decomposition and frequency-domain attention for long-horizon forecasting \citep{zhou2022fedformer}. PatchTST segments time series into patch tokens and uses channel-independent sharing for efficient long-context modeling \citep{nie2023patchtst}. TimesNet maps one-dimensional temporal variation into two-dimensional periodic structures and reports broad performance across forecasting, imputation, classification, and anomaly detection \citep{wu2023timesnet}. iTransformer inverts the standard representation and treats each variate as a token, allowing attention to operate over variables rather than individual timestamps \citep{liu2024itransformer}.

The variate-token choice is natural for AWS imputation because the core physical couplings are cross-variable: temperature, pressure, humidity, and wind are not independent channels. ECBIT adopts the iTransformer-style variate-token view, but it changes the task from forecasting to reconstruction under sparse observation masks. Each token contains a full 168-step trajectory plus mask and time features. This allows the model to learn variable-level relationships without treating every timestamp as a separate attention token.

\subsection{Exogenous Conditioning and Domain-Specific Restoration}
Recent work on auxiliary information in time series cautions against unconstrained fusion. \citet{lee2026constrainedfusion} report that naive multimodal fusion can underperform when auxiliary modalities inject irrelevant information, and propose the Controlled Fusion Adapter (CFA) for time-series forecasting settings. \citet{cao2026ecto} motivate physically grounded selection of meteorological exogenous variables rather than treating all auxiliary channels as uniformly useful. \citet{chen2025exost} identify inconsistent effects across exogenous variables and propose a ``select, then balance'' ExoST framework that dynamically selects and balances exogenous signals.

Meteorological restoration work provides a closer domain analogue. ST-DAN combines a transformer temporal path with a graph-attention path to restore buoy meteorological data using ERA5 and in situ observations \citep{song2026stdan}. ECBIT shares the motivation that physical variables and external reanalysis can help data restoration, but differs in three respects. First, it targets Antarctic AWS outages rather than marine buoy records. Second, it explicitly evaluates block-missing training against MCAR training, showing that the missingness curriculum is itself a major factor. Third, the ablation results show that a lightweight conditioning mechanism is sufficient once ERA5 alignment and block-curriculum matching are correct. The empirical value therefore lies less in a complex fusion block and more in matching the physical failure mode while controlling how external meteorological context enters the model.
